\documentclass{article}
\usepackage{amsmath, amssymb, amsthm}

\newtheorem{definition}{Definition}
\newtheorem{theorem}{Theorem}
\newtheorem{lemma}[theorem]{Lemma}

\begin{document}

\section*{The Fundamental Theorem of Finitely Generated Abelian Groups}

\begin{definition}[Group]
A \emph{group} is a set $G$ equipped with a binary operation $\cdot : G \times G \to G$ satisfying:
\begin{enumerate}
    \item (Associativity) $(a \cdot b) \cdot c = a \cdot (b \cdot c)$ for all $a,b,c \in G$.
    \item (Identity) There exists $e \in G$ with $e \cdot a = a \cdot e = a$ for all $a \in G$.
    \item (Inverses) For every $a \in G$ there exists $a^{-1} \in G$ with $a \cdot a^{-1} = a^{-1} \cdot a = e$.
\end{enumerate}
\end{definition}

\begin{definition}[Abelian group]
A group $(G, \cdot)$ is \emph{abelian} if $a \cdot b = b \cdot a$ for all $a, b \in G$. We typically write the operation additively: $a + b$, with identity $0$ and inverse $-a$.
\end{definition}

\begin{definition}[Subgroup and homomorphism]
A subset $H \subseteq G$ is a \emph{subgroup} if it is closed under the operation and inverses and contains $0$. A \emph{homomorphism} $\varphi : G \to H$ of abelian groups is a map satisfying $\varphi(a+b) = \varphi(a) + \varphi(b)$. An \emph{isomorphism} is a bijective homomorphism, written $G \cong H$.
\end{definition}

\begin{definition}[Direct sum]
Given abelian groups $A_1, \dots, A_n$, their \emph{direct sum} is
\[
    A_1 \oplus \cdots \oplus A_n 
    = \{ (a_1, \dots, a_n) : a_i \in A_i \},
\]
with componentwise addition $(a_i) + (b_i) = (a_i + b_i)$.
\end{definition}

\begin{definition}[Cyclic group]
An abelian group $G$ is \emph{cyclic} if there exists $g \in G$ such that $G = \{ ng : n \in \mathbb{Z} \}$. Every cyclic group is isomorphic to either $\mathbb{Z}$ (infinite order) or $\mathbb{Z}/n\mathbb{Z}$ for some $n \ge 1$ (finite order $n$).
\end{definition}

\begin{definition}[Finitely generated]
An abelian group $G$ is \emph{finitely generated} if there exist $g_1, \dots, g_n \in G$ such that every element can be written as $k_1 g_1 + \cdots + k_n g_n$ with $k_i \in \mathbb{Z}$.
\end{definition}

\begin{definition}[Torsion and free part]
An element $g \in G$ is a \emph{torsion element} if $ng = 0$ for some $n \ge 1$. The set of torsion elements forms a subgroup $T(G) \subseteq G$, called the \emph{torsion subgroup}. A group is \emph{torsion-free} if $T(G) = 0$, and \emph{free abelian of rank $r$} if it is isomorphic to $\mathbb{Z}^r$.
\end{definition}

\begin{theorem}[Fundamental Theorem, Invariant Factor Form]
Let $G$ be a finitely generated abelian group. Then there exist a unique integer $r \ge 0$ and unique integers $1 < d_1 \mid d_2 \mid \cdots \mid d_k$ such that
\[
    G \;\cong\; \mathbb{Z}^r \;\oplus\; \mathbb{Z}/d_1\mathbb{Z} \;\oplus\; \mathbb{Z}/d_2\mathbb{Z} \;\oplus\; \cdots \;\oplus\; \mathbb{Z}/d_k\mathbb{Z}.
\]
The integer $r$ is the \emph{rank} (or \emph{Betti number}) of $G$, and $d_1, \dots, d_k$ are the \emph{invariant factors}.
\end{theorem}

\begin{theorem}[Fundamental Theorem, Elementary Divisor Form]
Equivalently, there exist primes $p_1, \dots, p_s$ (not necessarily distinct) and positive integers $e_1, \dots, e_s$ such that
\[
    G \;\cong\; \mathbb{Z}^r \;\oplus\; \bigoplus_{i=1}^{s} \mathbb{Z}/p_i^{e_i}\mathbb{Z},
\]
and the multiset of prime powers $\{p_i^{e_i}\}$ is uniquely determined by $G$. These are the \emph{elementary divisors} of $G$.
\end{theorem}

\begin{lemma}[Passage between the two forms]
The two decompositions are related by the Chinese Remainder Theorem: if $d = p_1^{a_1} \cdots p_m^{a_m}$ is a prime factorization, then
\[
    \mathbb{Z}/d\mathbb{Z} \;\cong\; \mathbb{Z}/p_1^{a_1}\mathbb{Z} \;\oplus\; \cdots \;\oplus\; \mathbb{Z}/p_m^{a_m}\mathbb{Z}.
\]
Splitting each invariant factor yields the elementary divisors; conversely, grouping elementary divisors by taking, for each prime, the largest power into $d_k$, the next largest into $d_{k-1}$, and so on, recovers the invariant factors.
\end{lemma}

\end{document}